% Solutions of Book 4, chapter 22 (regular expressions and the limits of regular languages).
\section*{Chapter \ref{ch:b2:regular-languages} --- Regular Expressions and the Limits of Regular Languages}

\begin{solution}{exo:b2:regular-languages:1}
\omoutput{code/bachelor-2/22-regular-languages/outputs/exo_match.txt}{The seven words.}
$(ab \mid b)^*$ is the words in which every $a$ is immediately followed by a $b$, and the final $a?$ allows one more
$a$ at the end. So $\varepsilon$, $a$, $bab$, $abba$ and $babab$ match; $aa$ and $aab$ do not, since their first $a$
is followed by an $a$.
\end{solution}

\begin{solution}{exo:b2:regular-languages:2}
\omoutput{code/bachelor-2/22-regular-languages/outputs/exo_ident.txt}{The two patterns on six names.}
\code{[a-z0-9_]+} accepts \code{9lives} and \code{2}, which start with a digit. The pattern \code{[a-z_][a-z0-9_]*}
asks for one letter or underscore, then any number of letters, digits or underscores: it rejects them, and still
rejects the empty name.
\end{solution}

\begin{solution}{exo:b2:regular-languages:3}
\omoutput{code/bachelor-2/22-regular-languages/outputs/exo_decimal.txt}{The expression against the hook's automaton.}
The expression is \code{[+-]?[0-9]+(.[0-9]+)?([eE][+-]?[0-9]+)?}: an optional sign, at least one digit, an optional
group of a point and at least one digit, and an optional exponent with its own optional sign. It agrees with the
automaton on all 137\,257 strings of length at most 6 over the seven characters, 1\,854 of which are numbers. A check
on a finite set is not a proof; the proof is \cref{exo:b2:regular-languages:9}'s checker, run on the two automata.
\end{solution}

\begin{solution}{exo:b2:regular-languages:4}
\omoutput{code/bachelor-2/22-regular-languages/outputs/exo_count.txt}{$c_0$ to $c_{12}$, from the denotation.}
A nonempty word of the language starts either with $b$, followed by any word of the language of length $n - 1$, or
with $ab$, followed by one of length $n - 2$; the two cases do not overlap, since the first letter tells them apart.
So $c_n = c_{n-1} + c_{n-2}$ with $c_0 = c_1 = 1$: the Fibonacci numbers.
\end{solution}

\begin{solution}{exo:b2:regular-languages:5}
\omoutput{code/bachelor-2/22-regular-languages/outputs/exo_follow.txt}{The sets of $a_1(b_2a_3 \mid b_4)^*a_5$.}
The star's body has $\mathrm{First} = \{2, 4\}$ and $\mathrm{Last} = \{3, 4\}$; the star adds $\{2, 4\}$ to the follow
sets of 3 and 4. The concatenation of $a_1$ with the star (nullable) gives $\mathrm{Follow}(1) = \{2, 4\}$, and the
concatenation with $a_5$ adds 5 to the follow sets of 1, 3 and 4. Inside $b_2a_3$, $\mathrm{Follow}(2) = \{3\}$. The
automaton has states 0 to 5: $0 \xrightarrow{a} 1$; from 1, 3 and 4: $b$ to 2, $b$ to 4 and $a$ to 5; $2
\xrightarrow{a} 3$; state 5 accepts.
\end{solution}

\begin{solution}{exo:b2:regular-languages:6}
\omoutput{code/bachelor-2/22-regular-languages/outputs/exo_min_bug.txt}{The classes found, and four words on the
merged automaton.}
States 1 and 2 both have one successor, in the accepting class, so the program merges them; it forgot that state 1
moves on $b$ and state 2 on $a$. The merged automaton accepts $aa$ and $bb$ as well as $ab$ and $ba$. The fix is to
complete the automaton first, with a sink, and to compare the successors letter by letter, as in
\cref{lst:b2:regular-languages:moore}: then 1 goes to the sink on $a$ and 2 goes to 3, and round 1 separates them.
\end{solution}

\begin{solution}{exo:b2:regular-languages:7}
\omoutput{code/bachelor-2/22-regular-languages/outputs/exo_eliminate.txt}{The expression found by \code{eliminate}, and
its check.}
State 0 (even) is the start and accepts; $r_{00} = b$, $r_{01} = a$, $r_{10} = a$, $r_{11} = b$, with $S
\xrightarrow{\varepsilon} 0$ and $0 \xrightarrow{\varepsilon} E$. Removing 0 gives $r_{SE} = b^*$, $r_{S1} = b^*a$,
$r_{11} = b \mid ab^*a$ and $r_{1E} = ab^*$. Removing 1 gives $r_{SE} = b^* \mid b^*a(b \mid ab^*a)^*ab^*$: no $a$,
or a first $a$, blocks with an even number of $a$, and a last $a$.
\end{solution}

\begin{solution}{exo:b2:regular-languages:8}
Let $L = \{a^nb^m : n \le m\}$. Pumping: if a DFA with $p$ states recognised $L$, then $w = a^pb^p \in L$ would split
as $xyz$ with $|xy| \le p$, so $y = a^j$ with $j \ge 1$, and $xy^2z = a^{p+j}b^p$ would be in $L$; it is not, since
$p + j > p$. \omterm{def:b2:regular-languages:residual}{Residuals}: for $i < j$, $b^i$ is in $(a^i)^{-1}L$ but not in $(a^j)^{-1}L$, so the words $a^0, a^1,
a^2, \dots$ have pairwise different \omterm{def:b2:regular-languages:residual}{residuals}, infinitely many, and \cref{thm:b2:regular-languages:myhill-nerode}
concludes.
\end{solution}

\begin{solution}{exo:b2:regular-languages:9}
\omcode{code/bachelor-2/22-regular-languages/listings/dfa.rs}{94}{113}{A breadth-first search over pairs of states}
The search starts from the pair of start states and follows each letter in both automata at once: this is the \omterm{def:b2:finite-automata:product}{product
automaton}, built only as far as it is reached. Each pair is entered once, with the word that first reached it; since a
breadth-first search reaches pairs in order of the length of that word, the first pair with different acceptance
gives a shortest word on which the languages differ. If no such pair is reached, the languages are equal
(\cref{prop:b2:regular-languages:decisions}). The cost is at most $n_A n_B |\Sigma|$ steps. \Cref{out:b2:regular-languages:equivalent}
shows it on five pairs of expressions.
\end{solution}

\begin{solution}{exo:b2:regular-languages:10}
If $u^{-1}L = v^{-1}L$, then $(ua)^{-1}L = a^{-1}(u^{-1}L) = a^{-1}(v^{-1}L) = (va)^{-1}L$, since $x \in (ua)^{-1}L$
iff $uax \in L$ iff $ax \in u^{-1}L$. So $\sim_L$ is a right congruence. By definition, the classes of $\sim_L$ are in
bijection with the \omterm{def:b2:regular-languages:residual}{residuals}, by $[u] \mapsto u^{-1}L$. By \cref{cor:b2:regular-languages:unique-minimal}, the \omterm{def:b2:regular-languages:minimal}{minimal
automaton} is the \omterm{def:b2:regular-languages:residual}{residual} automaton, whose run on $u$ ends in $u^{-1}L$
(\cref{thm:b2:regular-languages:myhill-nerode}). So $u$ and $v$ lead it to the same state exactly when $u^{-1}L =
v^{-1}L$, that is, when $u \sim_L v$, and the number of classes, the \emph{index} of $\sim_L$, is the number of states
of the \omterm{def:b2:regular-languages:minimal}{minimal automaton}.
\end{solution}

\begin{solution}{exo:b2:regular-languages:11}
\omcode{code/bachelor-2/22-regular-languages/exercises/hopcroft.rs}{21}{52}{Hopcroft's loop: splitters, the states
entering them, and the smaller half}
\omoutput{code/bachelor-2/22-regular-languages/outputs/minimise_bench.txt}{Moore and Hopcroft on automata of up to
49\,153 states, on the laptop.}
The predecessor lists \code{inv} let a splitter $(B, a)$ cost the number of transitions entering $B$ on $a$. A state
$q$ lies in a splitter on the list once per time its class is put there, and after its first time it is put there
only as the smaller half, at most half of the previous class: at most $1 + \log_2 n$ times per letter. Summed over the
states, the work is proportional to $n|\Sigma| \log n$. The tests check that the classes are Moore's on every sample.
Here Hopcroft's algorithm gains less than a factor of 2 (between 1.25 and 1.6 in the run printed): Moore's method needs only $k + 1$ rounds on these automata, about $\log_2 n$, so both do about $n \log
n$ work, and Hopcroft's bookkeeping (a set of splitters, predecessor lists, sorting) costs more per step. On an automaton that is a chain of $n$ states, Moore needs about $n$ rounds and the
difference becomes large.
\end{solution}

\begin{solution}{exo:b2:regular-languages:12}
Let $L$ be the words with as many $a$ as $b$. \omterm{def:b2:regular-languages:residual}{Residuals}: for $i \ne j$, $b^i \in (a^i)^{-1}L$ and $b^i \notin
(a^j)^{-1}L$, so $L$ has infinitely many \omterm{def:b2:regular-languages:residual}{residuals}. Intersection: $L \cap a^*b^* = \{a^nb^n\}$, which is not \omterm{def:b2:regular-languages:regular}{regular}
(\cref{prop:b2:regular-languages:nonregular}); since $a^*b^*$ is \omterm{def:b2:regular-languages:regular}{regular} and \omterm{def:b2:regular-languages:regular}{regular languages} are closed under
intersection, $L$ cannot be \omterm{def:b2:regular-languages:regular}{regular}.
\end{solution}

\begin{solution}{pb:b2:regular-languages:1}
\textbf{1.} 4 end in Road, 5 in Row and 5 in Way: the 14 of the second line.

\textbf{2.} \code{[A-Z]} is an alternation of 26 letters, \code{[a-z]*} the star of 26 more, then the space and the
four letters of Road: $26 + 26 + 1 + 4 = 57$.

\textbf{3.} \code{[A-Z][a-z]* (Lane|Avenue)}: the 6 lanes and the 3 avenues, 9 names.

\textbf{4.} Bridge Road, Bridge Way, Castle Road, Castle Row, Castle Way, Garden Lane, Garden Road, Station Row.

\textbf{5.} The alternation of the names, each written as a concatenation of its letters, denotes the set.

\textbf{6.} $\Sigma^* e \Sigma^*$, with $\Sigma$ written as a class of every character that can occur in a name.

\textbf{7.} $a_1 b_2 b_3 c_4$; $\mathrm{Null}$ is false, $\mathrm{First} = \{1, 3, 4\}$, $\mathrm{Last} = \{4\}$,
$\mathrm{Follow}(1) = \{2\}$, $\mathrm{Follow}(2) = \mathrm{Follow}(3) = \{1, 3, 4\}$, $\mathrm{Follow}(4) =
\emptyset$.

\textbf{8.} 5 states, and $3 + 1 + 3 + 3 + 0 = 10$ transitions.

\textbf{9.} $\{0\} \xrightarrow{b} \{3\} \xrightarrow{a} \{1\} \xrightarrow{b} \{2\} \xrightarrow{b} \{3\}
\xrightarrow{c} \{4\}$: accepted, after 5 unions, one per letter, since each set holds one state.

\textbf{10.} The next set is masked by the positions holding the letter just read; equivalently, every transition into
a position carries that position's letter.

\textbf{11.} At most $(57 + 1) \times 12 = 696$ (\cref{prop:b2:regular-languages:linear}).

\textbf{12.} A set of states is the bits of one 128-bit integer, state 0 included. With a vector of 64-bit words as the
set, a union costs one OR per word, $\lceil (m + 1)/64 \rceil$ in all.

\textbf{13.} Six states: $\{0\}$; $\{1\}$, $\{3\}$ and $\{4\}$, reached by $a$, $b$ and $c$; $\{2\}$, reached from
$\{1\}$ by $b$; and the sink $\emptyset$. $\{2\}$ and $\{3\}$ have the transitions of $\{0\}$.

\textbf{14.} Round 0: $\{4\}$ against the rest. Round 1: $\{0\}$, $\{2\}$ and $\{3\}$ go to $\{4\}$ on $c$, the others
do not: three classes. Round 2: $\{1\}$ goes to $\{2\}$ on $b$, $\emptyset$ to $\emptyset$, which are now in different
classes: four classes. Round 3 splits nothing. The \omterm{def:b2:regular-languages:minimal}{minimal automaton} has 4 states.

\textbf{15.} $L$ (by $\varepsilon$, $b$, $ab$, \dots), $bL$ (by $a$), $\{\varepsilon\}$ (by $c$) and $\emptyset$ (by
$aa$, $cc$, \dots): the classes $\{\{0\}, \{2\}, \{3\}\}$, $\{1\}$, $\{4\}$ and $\emptyset$.

\textbf{16.} Both denote the words over $\{a, b\}$ in which every $a$ is immediately followed by a $b$, then $c$. Such
a word is a block of $b$'s, then blocks starting with $ab$ followed by $b$'s, which is $b^*(abb^*)^*$; and every word of
$b^*(abb^*)^*$ has each $a$ followed by $b$. The checker of \cref{exo:b2:regular-languages:9} confirms it.

\textbf{17.} $c$: it is in the first language and lacks the $b$ before $c$ in the second. No shorter word differs,
since $\varepsilon$, $a$ and $b$ are in neither.

\textbf{18.} By \cref{cor:b2:regular-languages:unique-minimal}, patterns with the same language have the same \omterm{def:b2:regular-languages:minimal}{minimal
automaton} up to the names of the states. Renumbering its states in the order a breadth-first search from the start
meets them, letters in a fixed order, makes the transition table and the \omterm{def:b2:finite-automata:dfa}{accepting states} a canonical key.

\textbf{19.} $2^n$: each $a$ is read by the left or the right branch.

\textbf{20.} Every way reaches the end of the text without a $b$ and fails; the engine then returns to its last choice
and tries the other branch, so it tries all $2^n$. One more $a$ doubles the ways, and the measured ratios are close
to 2.

\textbf{21.} About $6.6 \times 2^{13} \approx 54\,000$ s, 15 hours.

\textbf{22.} One union for the first letter, from $\{0\}$, then two for each later one, from $\{1, 2\}$: $2n - 1$, as
the output shows (39 for 20 letters, 1\,999 for 1\,000).

\textbf{23.} About 7.3 ns a letter: some 0.3 $\mu$s for 40 letters, and 73 ms for ten million.

\textbf{24.} \emph{From exponential to linear}: on $(a \mid a)^*b$ and $a^n$, the backtracking engine's time doubles
with each $a$, as measured, reaching 6.6 s at $n = 27$; the Glushkov matcher makes at most $(m + 1)n$ unions on any
pattern of $m$ positions and any text of $n$ letters, as proved, here $2n - 1$, and answers ten million letters in
73 ms.
\end{solution}
